% ═══════════════════════════════════════════════════════════════
% FOLIEN: parecer y conocer (DS 1)
% Present/Apply-Struktur für Stundenführung
% ═══════════════════════════════════════════════════════════════

\documentclass[aspectratio=169, 11pt]{beamer}

\usepackage[utf8]{inputenc}
\usepackage[T1]{fontenc}
\usepackage[ngerman]{babel}

% --- SCHRIFTEN ---
\usepackage{helvet}
\renewcommand{\familydefault}{\sfdefault}
\usepackage{palatino}
\newcommand{\seriftext}[1]{{\fontfamily{ppl}\selectfont #1}}

\usepackage{tikz}
\usepackage{array}
\usepackage{booktabs}
\usepackage{pifont}

% ═══════════════════════════════════════════════════════════════
% FARBEN
% ═══════════════════════════════════════════════════════════════

\definecolor{akzent}{HTML}{C41E3A}      % Karminrot (Spanisch)
\definecolor{stamm}{HTML}{C62828}       % Stammwechsel
\definecolor{hellgrau}{HTML}{F5F5F5}
\definecolor{dunkelgrau}{HTML}{333333}
\definecolor{mittelgrau}{HTML}{666666}
\definecolor{richtig}{HTML}{2A7A4B}     % Grün für Haken
\definecolor{falsch}{HTML}{C62828}      % Rot für Kreuz

% ═══════════════════════════════════════════════════════════════
% BEAMER-THEME
% ═══════════════════════════════════════════════════════════════

\usetheme{default}
\usecolortheme{default}

\setbeamertemplate{navigation symbols}{}

\setbeamercolor{frametitle}{fg=white, bg=akzent}
\setbeamercolor{title}{fg=white}
\setbeamercolor{subtitle}{fg=white}
\setbeamercolor{author}{fg=white}
\setbeamercolor{date}{fg=white}
\setbeamercolor{structure}{fg=akzent}
\setbeamercolor{itemize item}{fg=akzent}
\setbeamercolor{itemize subitem}{fg=mittelgrau}
\setbeamercolor{block title}{fg=white, bg=akzent}
\setbeamercolor{block body}{bg=hellgrau}

\setbeamertemplate{itemize item}{\textbullet}
\setbeamertemplate{itemize subitem}{--}

\setbeamertemplate{title page}{
    \begin{tikzpicture}[remember picture, overlay]
        \fill[akzent] (current page.north west) rectangle (current page.south east);
    \end{tikzpicture}
    \vspace{2cm}
    \begin{center}
        {\Huge\bfseries\textcolor{white}{\inserttitle}}\\[0.8cm]
        {\Large\textcolor{white}{\insertsubtitle}}\\[1.5cm]
        {\large\textcolor{white}{\insertauthor}}\\[0.3cm]
        {\textcolor{white}{\insertdate}}
    \end{center}
}

\setbeamertemplate{frametitle}{
    \nointerlineskip
    \begin{beamercolorbox}[wd=\paperwidth, ht=1.2cm, dp=0.3cm]{frametitle}
        \hspace{0.5cm}\Large\bfseries\insertframetitle
    \end{beamercolorbox}
}

\setbeamertemplate{footline}{
    \hfill\textcolor{mittelgrau}{\footnotesize\insertframenumber}\hspace{0.5cm}\vspace{0.3cm}
}

% ═══════════════════════════════════════════════════════════════
% BEFEHLE
% ═══════════════════════════════════════════════════════════════

\newcommand{\es}[1]{\seriftext{\textbf{#1}}}
\newcommand{\de}[1]{\textit{\textcolor{mittelgrau}{#1}}}
\newcommand{\stw}[1]{\textcolor{stamm}{\textbf{#1}}}
\newcommand{\abmark}[1]{{\Large$\rightarrow$ \textbf{Arbeitsblatt Aufgabe #1}}}

% ═══════════════════════════════════════════════════════════════
% METADATEN
% ═══════════════════════════════════════════════════════════════

\title{parecer y conocer}
\subtitle{Meinungen äußern und Personen kennen}
\author{Spanisch · Klasse 9}
\date{}

% ═══════════════════════════════════════════════════════════════
% DOKUMENT
% ═══════════════════════════════════════════════════════════════

\begin{document}

% ═══════════════════════════════════════════════════════════════
% EINSTIEG
% ═══════════════════════════════════════════════════════════════

\begin{frame}[plain]
    \titlepage
\end{frame}

\begin{frame}{Lernziele}
    \vspace{0.5cm}

    {\large Heute lernst du:}

    \vspace{0.8cm}

    \begin{itemize}
        \setlength{\itemsep}{0.6cm}
        \item \textbf{parecer} \de{-- scheinen, vorkommen}\\
              {\small Funktioniert wie \es{gustar}!}
        \item \textbf{conocer} \de{-- kennen, kennenlernen}\\
              {\small Mit Stammwechsel in der 1. Person}
        \item \textbf{conocer + a} bei Personen\\
              {\small Die persönliche Ergänzung mit \es{a}}
    \end{itemize}
\end{frame}

% ═══════════════════════════════════════════════════════════════
% BLOCK 1: VOCABULARIO
% ═══════════════════════════════════════════════════════════════

\begin{frame}{Vocabulario}
    \vspace{0.3cm}

    \begin{columns}[T]
        \begin{column}{0.45\textwidth}
            \textbf{Verben}

            \vspace{0.4cm}

            \begin{tabular}{@{}ll@{}}
                \es{parecer} & \de{scheinen} \\[0.4em]
                \es{conocer} & \de{kennen} \\
            \end{tabular}

            \vspace{0.6cm}

            \textbf{Adverbien}

            \vspace{0.4cm}

            \begin{tabular}{@{}ll@{}}
                \es{todavía no} & \de{noch nicht} \\[0.4em]
                \es{bastante} & \de{ziemlich} \\
            \end{tabular}
        \end{column}

        \begin{column}{0.45\textwidth}
            \textbf{Adjektive}

            \vspace{0.4cm}

            \begin{tabular}{@{}ll@{}}
                \es{simpático/-a} & \de{sympathisch} \\[0.4em]
                \es{engreído/-a} & \de{eingebildet} \\[0.4em]
                \es{tranquilo/-a} & \de{ruhig} \\[0.4em]
                \es{raro/-a} & \de{seltsam} \\[0.4em]
                \es{guapo/-a} & \de{gutaussehend} \\
            \end{tabular}
        \end{column}
    \end{columns}
\end{frame}

% ═══════════════════════════════════════════════════════════════
% BLOCK 2: PARECER (PRESENT)
% ═══════════════════════════════════════════════════════════════

\begin{frame}{parecer -- wie gustar!}
    \vspace{0.3cm}

    \begin{block}{Merke}
        \es{parecer} funktioniert genau wie \es{gustar}:\\
        \textbf{Indirektes Objektpronomen + parecer/parecen}
    \end{block}

    \vspace{0.6cm}

    \begin{columns}[T]
        \begin{column}{0.45\textwidth}
            \textbf{Singular} (eine Sache)

            \vspace{0.3cm}

            \es{Me \textbf{parece} bien.}\\
            \de{Ich finde es gut.}

            \vspace{0.3cm}

            \es{¿Qué te \textbf{parece}?}\\
            \de{Was meinst du?}
        \end{column}

        \begin{column}{0.45\textwidth}
            \textbf{Plural} (mehrere Sachen)

            \vspace{0.3cm}

            \es{Me \textbf{parecen} simpáticos.}\\
            \de{Ich finde sie sympathisch.}

            \vspace{0.3cm}

            \es{¿Qué te \textbf{parecen}?}\\
            \de{Was hältst du von ihnen?}
        \end{column}
    \end{columns}
\end{frame}

\begin{frame}{parecer -- Konjugation}
    \vspace{0.5cm}

    \begin{center}
        {\large\textbf{Indirektes Objektpronomen + parece(n)}}

        \vspace{0.8cm}

        \begin{tabular}{lll}
            \toprule
            \textbf{Pronomen} & \textbf{Singular} & \textbf{Plural} \\
            \midrule
            me & \es{me parece} & \es{me parecen} \\[0.4em]
            te & \es{te parece} & \es{te parecen} \\[0.4em]
            le & \es{le parece} & \es{le parecen} \\[0.4em]
            nos & \es{nos parece} & \es{nos parecen} \\[0.4em]
            os & \es{os parece} & \es{os parecen} \\[0.4em]
            les & \es{les parece} & \es{les parecen} \\
            \bottomrule
        \end{tabular}
    \end{center}
\end{frame}

% --- APPLY: parecer üben ---
\begin{frame}{Jetzt du! -- parecer}
    \vspace{1cm}

    \begin{center}
        \abmark{1}

        \vspace{1cm}

        {\large Setze die richtige Form von \es{parecer} ein.}

        \vspace{0.8cm}

        \begin{tabular}{cl}
            \textcolor{akzent}{\textbf{1.}} & Partnerarbeit \\[0.3em]
            \textcolor{akzent}{\textbf{2.}} & Vergleicht eure Lösungen \\
        \end{tabular}

        \vspace{0.8cm}

        {\small Zeit: 5 Minuten}
    \end{center}
\end{frame}

% ═══════════════════════════════════════════════════════════════
% BLOCK 3: CONOCER (PRESENT)
% ═══════════════════════════════════════════════════════════════

\begin{frame}{conocer -- Konjugation}
    \vspace{0.3cm}

    \begin{columns}[T]
        \begin{column}{0.45\textwidth}
            \begin{center}
                {\large\textbf{Konjugation}}

                \vspace{0.5cm}

                \begin{tabular}{ll}
                    yo & \es{cono\stw{zc}o} \\[0.4em]
                    tú & \es{conoces} \\[0.4em]
                    él/ella & \es{conoce} \\[0.4em]
                    nosotros & \es{conocemos} \\[0.4em]
                    vosotros & \es{conocéis} \\[0.4em]
                    ellos & \es{conocen} \\
                \end{tabular}
            \end{center}
        \end{column}

        \begin{column}{0.45\textwidth}
            \vspace{1cm}

            \begin{block}{Stammwechsel}
                \textbf{Nur 1. Person Singular!}

                \vspace{0.3cm}

                c → \stw{zc} vor -o

                \vspace{0.3cm}

                yo \es{cono\stw{zc}o}
            \end{block}

            \vspace{0.3cm}

            {\small Auch bei: \es{traducir, producir, conducir...}}
        \end{column}
    \end{columns}
\end{frame}

\begin{frame}{conocer + a bei Personen}
    \vspace{0.5cm}

    \begin{block}{Wichtige Regel}
        Bei \textbf{Personen} steht immer \es{\textbf{a}} nach \es{conocer}!
    \end{block}

    \vspace{0.6cm}

    \begin{columns}[T]
        \begin{column}{0.45\textwidth}
            {\large\textcolor{richtig}{\ding{51}} \textbf{Person} → mit \es{a}}

            \vspace{0.4cm}

            \es{¿Conoces \textbf{a} María?}\\
            \de{Kennst du María?}

            \vspace{0.3cm}

            \es{Conozco \textbf{a} tu hermano.}\\
            \de{Ich kenne deinen Bruder.}
        \end{column}

        \begin{column}{0.45\textwidth}
            {\large\textcolor{falsch}{\ding{55}} \textbf{Sache} → ohne a}

            \vspace{0.4cm}

            \es{¿Conoces la ciudad?}\\
            \de{Kennst du die Stadt?}

            \vspace{0.3cm}

            \es{Conozco este libro.}\\
            \de{Ich kenne dieses Buch.}
        \end{column}
    \end{columns}
\end{frame}

% --- APPLY: conocer üben ---
\begin{frame}{Jetzt du! -- conocer}
    \vspace{0.8cm}

    \begin{center}
        \abmark{2 + 3}

        \vspace{0.8cm}

        \begin{tabular}{cl}
            \textcolor{akzent}{\textbf{Aufgabe 2:}} & Konjugiere \es{conocer} \\[0.5em]
            \textcolor{akzent}{\textbf{Aufgabe 3:}} & Entscheide: mit \es{a} oder ohne? \\
        \end{tabular}

        \vspace{0.8cm}

        {\large Einzelarbeit, dann Partnerkontrolle}

        \vspace{0.5cm}

        {\small Zeit: 8 Minuten}
    \end{center}
\end{frame}

% ═══════════════════════════════════════════════════════════════
% BLOCK 4: KOMMUNIKATION (PRESENT)
% ═══════════════════════════════════════════════════════════════

\begin{frame}{Redemittel -- Meinung äußern}
    \vspace{0.3cm}

    \begin{columns}[T]
        \begin{column}{0.45\textwidth}
            \textbf{Fragen}

            \vspace{0.4cm}

            \es{¿Qué te parece...?}\\
            \de{Was meinst du zu...?}

            \vspace{0.3cm}

            \es{¿Qué te parecen...?}\\
            \de{Was hältst du von...?}

            \vspace{0.3cm}

            \es{¿Conoces a...?}\\
            \de{Kennst du...?}
        \end{column}

        \begin{column}{0.45\textwidth}
            \textbf{Antworten}

            \vspace{0.4cm}

            \es{Me parece bien/mal.}\\
            \de{Ich finde es gut/schlecht.}

            \vspace{0.3cm}

            \es{Me parece interesante.}\\
            \de{Ich finde es interessant.}

            \vspace{0.3cm}

            \es{Sí, lo/la conozco.}\\
            \de{Ja, ich kenne ihn/sie.}
        \end{column}
    \end{columns}
\end{frame}

\begin{frame}{Diálogo modelo}
    \vspace{0.3cm}

    \begin{center}
        \begin{tabular}{@{}p{7cm}p{5.5cm}@{}}
            \textbf{A:} \es{¿Conoces a mi hermano?} & \de{Kennst du meinen Bruder?} \\[0.5em]
            \textbf{B:} \es{No, todavía no lo conozco.} & \de{Nein, noch nicht.} \\[0.5em]
            \textbf{B:} \es{¿Cómo es?} & \de{Wie ist er so?} \\[0.5em]
            \textbf{A:} \es{Me parece un poco raro.} & \de{Ich finde ihn etwas seltsam.} \\[0.5em]
            \textbf{B:} \es{¿Qué te parece su música?} & \de{Was hältst du von seiner Musik?} \\[0.5em]
            \textbf{A:} \es{Me parece bastante buena.} & \de{Ziemlich gut, finde ich.} \\
        \end{tabular}
    \end{center}
\end{frame}

% --- APPLY: Dialog schreiben ---
\begin{frame}{Jetzt du! -- Dialog}
    \vspace{0.8cm}

    \begin{center}
        \abmark{4}

        \vspace{0.8cm}

        {\large Schreibt einen eigenen Dialog mit \es{parecer} und \es{conocer}.}

        \vspace{0.6cm}

        \begin{tabular}{cl}
            \textcolor{akzent}{\textbf{1.}} & Partnerarbeit (6 Min.) \\[0.3em]
            \textcolor{akzent}{\textbf{2.}} & Dialog vorlesen (freiwillig) \\
        \end{tabular}

        \vspace{0.8cm}

        {\small Nutzt die Redemittel und den Beispieldialog!}
    \end{center}
\end{frame}

% ═══════════════════════════════════════════════════════════════
% SICHERUNG
% ═══════════════════════════════════════════════════════════════

\begin{frame}{Zusammenfassung}
    \vspace{0.3cm}

    \begin{columns}[T]
        \begin{column}{0.45\textwidth}
            \textbf{1. parecer}

            \vspace{0.3cm}

            \begin{itemize}
                \setlength{\itemsep}{0.3cm}
                \item Wie \es{gustar}
                \item Ind. Objektpronomen + \es{parece(n)}
                \item \es{parece} = Singular
                \item \es{parecen} = Plural
            \end{itemize}
        \end{column}

        \begin{column}{0.45\textwidth}
            \textbf{2. conocer}

            \vspace{0.3cm}

            \begin{itemize}
                \setlength{\itemsep}{0.3cm}
                \item yo = \es{cono\stw{zc}o}
                \item Stammwechsel c → zc
                \item \textbf{+ a} bei Personen!
                \item Ohne a bei Sachen
            \end{itemize}
        \end{column}
    \end{columns}
\end{frame}

\begin{frame}{Typische Fehler vermeiden}
    \vspace{0.5cm}

    \begin{center}
        \begin{tabular}{@{}cl@{\hspace{2cm}}cl@{}}
            \textcolor{falsch}{\Large$\times$} & \es{Conozco María.} & \textcolor{richtig}{\Large$\checkmark$} & \es{Conozco \textbf{a} María.} \\[0.8em]
            \textcolor{falsch}{\Large$\times$} & \es{Yo conoco...} & \textcolor{richtig}{\Large$\checkmark$} & \es{Yo cono\stw{zc}o...} \\[0.8em]
            \textcolor{falsch}{\Large$\times$} & \es{Me parece simpáticos.} & \textcolor{richtig}{\Large$\checkmark$} & \es{Me \textbf{parecen} simpáticos.} \\[0.8em]
            \textcolor{falsch}{\Large$\times$} & \es{Te parece las clases?} & \textcolor{richtig}{\Large$\checkmark$} & \es{¿Te \textbf{parecen} las clases?} \\
        \end{tabular}
    \end{center}
\end{frame}

\begin{frame}{Geschafft!}
    \vspace{1cm}

    \begin{center}
        {\Large\textcolor{richtig}{\ding{51}} Du kannst jetzt:}

        \vspace{0.8cm}

        \begin{itemize}
            \setlength{\itemsep}{0.5cm}
            \item Mit \es{parecer} deine Meinung sagen
            \item \es{conocer} richtig konjugieren
            \item \es{a} bei Personen nicht vergessen
        \end{itemize}

        \vspace{1cm}

        {\large\textcolor{akzent}{\textbf{¡Muy bien!}}}
    \end{center}
\end{frame}

\end{document}
